\documentclass[a4paper,12pt]{article}
\usepackage[T2A]{fontenc}
\usepackage[utf8]{inputenc}
\usepackage[russian]{babel}
\usepackage{amsmath,amssymb,mathtools,icomma}
\usepackage[a4paper,left=20mm,right=20mm,top=23mm,bottom=19mm,headheight=15pt]{geometry}
\usepackage{microtype,parskip,setspace,fancyhdr,xcolor,enumitem,booktabs,tabularx,titlesec}
\usepackage{tikz,pgfplots}
\usetikzlibrary{arrows.meta,positioning,shapes.geometric,babel}
\pgfplotsset{compat=1.18}
\definecolor{accent}{HTML}{008080}
\definecolor{darkaccent}{HTML}{004D4D}
\definecolor{textgray}{HTML}{333333}
\color{textgray}
\setstretch{1.25}
\setlength{\parskip}{6pt}
\setlength{\parindent}{0pt}
\setlist[itemize]{leftmargin=*,itemsep=1pt,topsep=3pt}
\setlist[enumerate]{leftmargin=*,itemsep=4pt,topsep=3pt}
\titleformat{\section}{\Large\normalfont\color{darkaccent}}{\thesection.}{0.5em}{}[\vspace{2pt}\color{accent}\titlerule]
\titlespacing*{\section}{0pt}{4pt}{12pt}
\pagestyle{fancy}
\fancyhf{}
\fancyhead[L]{\small Корни, степени и уравнения}
\fancyhead[R]{\small Ксения Клыкова}
\fancyfoot[C]{\small\thepage}
\renewcommand{\headrulewidth}{0.4pt}
\newcommand{\topic}[1]{\par\medskip{\color{darkaccent}\textbf{#1}}\par\smallskip}
\newcommand{\checkitem}{\(\square\)\enspace}
\begin{document}
\begin{titlepage}
\centering
\vspace*{27mm}
{\Large\color{darkaccent} МАТЕМАТИКА\par}
\vspace{15mm}
{\Huge Чек-лист\par}
\vspace{10mm}
{\LARGE\color{accent} Корни и дробные степени.\\ Область определения и\\ иррациональные уравнения\par}
\vspace{14mm}
{\large Алгебра. Подготовка к контрольной работе\par}
\vfill
{\large Ксения Клыкова\par}
\vspace{9mm}
{\large 26 ноября 2026 г.\par}
\vspace{16mm}
{Лёвин Артём Александрович\\эксклюзивно для Ксении Клыковой\par}
\vspace{20mm}
\begin{tikzpicture}
\draw[accent,line width=.9pt] (0,0) .. controls (2,0.75) and (4,-0.75) .. (6,0) .. controls (8,0.75) and (10,-0.75) .. (12,0);
\end{tikzpicture}
\vspace{7mm}
\end{titlepage}
\section{Область определения функции}
Область определения --- множество всех значений аргумента, при которых формула функции имеет смысл в действительных числах. ОДЗ --- область допустимых значений переменной.

\begin{center}
\renewcommand{\arraystretch}{1.5}
\begin{tabularx}{\linewidth}{@{}lXl@{}}
\toprule
Выражение & Условие & Почему \\
\midrule
$\sqrt{F(x)}$ & $F(x)\ge 0$ & Чётный корень \\
$\dfrac{1}{G(x)}$ & $G(x)\ne 0$ & Деление на нуль \\
$\dfrac{1}{\sqrt{F(x)}}$ & $F(x)>0$ & Оба ограничения \\
\bottomrule
\end{tabularx}
\end{center}
\topic{Пример 1. Корень в знаменателе}
Найдите область определения функции $y=\dfrac{1}{\sqrt{2x-x^2}}$.
\[
\begin{cases}2x-x^2\ge0,\\2x-x^2\ne0\end{cases}
\iff 2x-x^2>0\iff x(2-x)>0.
\]
Нули множителей: $x=0$ и $x=2$. На среднем промежутке проверяем $x=1$: $1\cdot(2-1)=1>0$. На внешних промежутках произведение отрицательно.
\begin{center}
\begin{tikzpicture}[x=1.65cm,y=1cm,>=Stealth]
\draw[-{Stealth}] (-1.6,0)--(4.0,0) node[right]{$x$};
\draw[accent,line width=2pt] (0,0)--(2,0);
\foreach \x in {0,2}{\node[circle,draw=accent,fill=white,line width=1pt,inner sep=0pt,minimum size=5pt] at (\x,0) {};\node[below=5pt,black] at (\x,0) {$\x$};}
\node at (-.8,.4){$-$};\node[accent] at (1,.4){$+$};\node at (3,.4){$-$};
\end{tikzpicture}
\end{center}
\[\boxed{D(y)=(0;2).}\]
\topic{Граничные точки и скобки}
\begin{itemize}
\item Строгое неравенство: нули выражения исключаются; пустая точка и круглая скобка.
\item Нестрогое неравенство: допустимый нуль включается; закрашенная точка и квадратная скобка.
\item Точка, в которой знаменатель равен нулю, исключается при любом знаке неравенства. У бесконечности всегда круглая скобка.
\end{itemize}
\clearpage
\section{Восстановление формулы по графику}
Если известен вид функции $f(x)=a\sqrt{x}$, достаточно одной точки графика с ненулевой абсциссой, чтобы определить коэффициент $a$.

Для точки $A(x_0;y_0)$ с $x_0>0$ выполняется
\[y_0=a\sqrt{x_0}\quad\Longrightarrow\quad a=\frac{y_0}{\sqrt{x_0}}.\]
\topic{Пример 2. Найдите $f(25)$}
На графике отмечена точка $A(1;4)$. Первая координата --- аргумент $x$, вторая --- значение функции $y$.
\begin{center}
\begin{tikzpicture}
\begin{axis}[width=7cm,height=10cm,axis lines=middle,xmin=-.6,xmax=5.5,ymin=-.7,ymax=9.8,xtick={1,2,3,4,5},ytick={2,4,6,8},xlabel={$x$},ylabel={$y$},grid=major,grid style={gray!18},clip=false,samples=220,axis equal image]
\addplot[accent,very thick,domain=0:5.3] {4*sqrt(x)};
\addplot[only marks,mark=*,mark size=2pt,darkaccent] coordinates {(1,4)};
\draw[dashed,gray] (axis cs:1,0)--(axis cs:1,4)--(axis cs:0,4);
\node[anchor=west,fill=white,inner sep=2pt] at (axis cs:1.15,3.8) {$A(1;4)$};
\end{axis}
\end{tikzpicture}
\end{center}
\begin{enumerate}
\item Подставляем координаты точки: $4=a\sqrt{1}$.
\item Находим коэффициент: $a=4$.
\item Восстанавливаем формулу: $f(x)=4\sqrt{x}$.
\item Подставляем требуемый аргумент: $f(25)=4\sqrt{25}=20$.
\end{enumerate}
\topic{Проверяйте}
Точка $(0;0)$ не позволяет определить $a$: равенство $0=a\cdot0$ верно при любом $a$. Значение $f(25)$ можно вычислить по восстановленной формуле, даже если соответствующая точка лежит за пределами рисунка.
\clearpage
\section{Дробные показатели степени}
При преобразованиях с дробными и отрицательными показателями удобно сначала работать с положительными основаниями. Для $a,b>0$ и рациональных $r,s$:
\[\sqrt[n]{a^m}=a^{m/n},\qquad a^{-r}=\frac1{a^r},\qquad a^ra^s=a^{r+s},\]
\[\frac{a^r}{a^s}=a^{r-s},\qquad (a^r)^s=a^{rs},\qquad (ab)^s=a^sb^s.\]
Корень чётной степени допускает также нулевое подкоренное выражение, если отсутствует деление на нуль.
\topic{Пример 3. Сложное степенное выражение}
При $a,b>0$ упростите
\[\left(b^{1\frac47}\cdot\sqrt{a^3}:a^{-\frac1{14}}\right)^{2\frac6{11}}.\]
Сначала заменяем корень степенью. При делении показатели вычитаются; отрицательный показатель записываем в скобках:
\[\left(b^{1\frac47}a^{\frac32-(-\frac1{14})}\right)^{2\frac6{11}},\qquad
\frac32+\frac1{14}=\frac{21+1}{14}=\frac{11}{7}.\]
Преобразуем смешанные числа: $1\frac47=\frac{11}{7}$, $2\frac6{11}=\frac{28}{11}$.
\begin{align*}
\left(b^{\frac{11}{7}}a^{\frac{11}{7}}\right)^{\frac{28}{11}}
&=b^{\frac{11}{7}\cdot\frac{28}{11}}a^{\frac{11}{7}\cdot\frac{28}{11}}\\
&=b^4a^4=(ab)^4.
\end{align*}
\topic{Вычислительная дисциплина}
\[\frac72-\frac78=\frac{28}{8}-\frac78=\frac{21}{8}.\]
\begin{itemize}
\item Перепишите все знаки, включая минус в показателе.
\item При сложении и вычитании дробей сначала получите общий знаменатель.
\item Перед умножением переведите смешанные числа в неправильные дроби и сократите множители.
\item При возведении произведения в степень примените показатель к каждому множителю.
\end{itemize}
\clearpage
\section{Сокращение дробей со степенями}
Чтобы сократить алгебраическую дробь, разложите её числитель и знаменатель на множители. Сокращать можно общие ненулевые множители. Ограничения исходного выражения сохраняются.
\topic{Два ключевых преобразования}
\[x^2y-x=x(xy-1),\qquad u^2-v^2=(u-v)(u+v).\]
При вынесении множителя проверяйте результат обратным раскрытием скобок.
\topic{Пример 4. Сократите дробь}
\[\frac{a^{1/3}-3a^{1/6}}{a^{1/3}-9}.\]
Здесь $a\ge0$, поскольку в выражении есть корень шестой степени. Знаменатель не равен нулю:
\[a^{1/3}\ne9\quad\Longrightarrow\quad a\ne729.\]
Представим $a^{1/3}=a^{2/6}=(a^{1/6})^2$.

\textbf{Числитель.} Выносим общий множитель:
\[a^{2/6}-3a^{1/6}=a^{1/6}(a^{1/6}-3).\]
Равенство верно и при $a=0$, что видно при раскрытии скобок.

\textbf{Знаменатель.} Используем разность квадратов:
\[a^{2/6}-9=(a^{1/6}-3)(a^{1/6}+3).\]
\textbf{Сокращение.} Множитель $a^{1/6}-3\ne0$ на исходной ОДЗ, поэтому
\[\frac{a^{1/6}(a^{1/6}-3)}{(a^{1/6}-3)(a^{1/6}+3)}
=\frac{a^{1/6}}{a^{1/6}+3}.\]
\[\boxed{\frac{a^{1/6}}{a^{1/6}+3},\qquad a\ge0,\quad a\ne729.}\]
\topic{Типичная ошибка}
При $a=729$ упрощённая дробь определена, однако исходная содержит деление на нуль. Поэтому исключённое значение нельзя возвращать после сокращения.
\clearpage
\section{Иррациональные уравнения}
\begingroup
\setstretch{1.15}
\setlength{\parskip}{3pt}
Иррациональное уравнение содержит неизвестную под знаком корня. Для квадратного корня ключевой равносильный переход:
\[\boxed{\sqrt{F(x)}=G(x)\iff
\begin{cases}G(x)\ge0,\\F(x)=G^2(x).\end{cases}}\]
При необходимости сохраняйте ограничения области определения самих выражений $F$ и $G$. Условие $F(x)\ge0$ уже следует из второго равенства. Первое условие обеспечивает правильный знак правой части, поскольку квадратный корень всегда неотрицателен.
\topic{Пример 5. Справа положительное число}
\[\sqrt{5-4x}=3,2\iff 5-4x=10,24\iff x=-1,31.\]
Условие $3,2\ge0$ выполнено. Проверка: $5-4(-1,31)=10,24$; корень равен $3,2$.
\topic{Пример 6. Справа переменная}
\[\sqrt{6+5x}=x\iff
\begin{cases}x\ge0,\\x^2-5x-6=0.\end{cases}\]
$(x-6)(x+1)=0$, откуда $x=6$ или $x=-1$. По условию $x\ge0$ подходит только $6$. Ответ: $6$.
\topic{Пример 7. Справа отрицание переменной}
\[\sqrt{-72-17x}=-x\iff
\begin{cases}x\le0,\\x^2+17x+72=0.\end{cases}\]
Здесь $(-x)^2=x^2$. $(x+9)(x+8)=0$. Оба корня $-9$ и $-8$ удовлетворяют $x\le0$. Ответ: $-9; -8$.
\topic{Пример 8. Справа сумма}
\[\sqrt{4x^2-3x-1}=x+1\iff
\begin{cases}x\ge-1,\\4x^2-3x-1=(x+1)^2.\end{cases}\]
Раскрываем квадрат: $(x+1)^2=x^2+2x+1$. Получаем
\[3x^2-5x-2=(3x+1)(x-2)=0.\]
Корни $-\frac13$ и $2$ удовлетворяют $x\ge-1$. Ответ: $-\frac13;2$.
\endgroup
\clearpage
\section*{Алгоритм: уравнение с квадратным корнем}
\vfill
\begin{center}
\begin{tikzpicture}[
node distance=10mm and 12mm,
start/.style={ellipse,draw=accent,minimum width=38mm,minimum height=10mm,align=center},
actionbox/.style={rectangle,rounded corners=2pt,draw=accent,text width=80mm,minimum height=12mm,align=center,inner sep=5pt},
choice/.style={diamond,aspect=2.5,draw=darkaccent,align=center,inner sep=2pt,text width=42mm},
arrow/.style={-{Stealth[length=2.5mm]},semithick,darkaccent}]
\node[start] (s) {$\sqrt{F(x)}=G(x)$};
\node[actionbox,below=of s] (d) {Учесть исходные ограничения выражений $F$ и $G$};
\node[actionbox,below=of d] (r) {Записать систему: $G(x)\ge0$ и $F(x)=[G(x)]^2$};
\node[actionbox,below=of r] (q) {Решить уравнение после возведения в квадрат};
\node[choice,below=of q] (c) {Есть кандидаты, удовлетворяющие всем ограничениям?};
\node[actionbox,below=15mm of c,text width=65mm] (yes) {Записать все допустимые корни};
\node[actionbox,right=of c,text width=29mm] (no) {Решений нет: $\emptyset$};
\node[start,below=of yes] (e) {Ответ};
\draw[arrow] (s)--(d);\draw[arrow](d)--(r);\draw[arrow](r)--(q);\draw[arrow](q)--(c);
\draw[arrow](c)--node[left]{да}(yes);\draw[arrow](c)--node[above]{нет}(no);
\draw[arrow](yes)--(e);\draw[arrow](no.south)|-(e.east);
\end{tikzpicture}
\end{center}
\vfill
\clearpage
\section*{Тренировочный блок}
Выполняйте преобразования последовательно. Указывайте ОДЗ, записывайте системы при решении иррациональных уравнений и проверяйте допустимость полученных корней.
\topic{Базовый уровень}
\begin{enumerate}[label=\arabic*.]
\item Найдите область определения функции $y=\sqrt{5-x}$.
\item Найдите область определения функции $y=\dfrac{1}{\sqrt{x(6-x)}}$.
\item График функции $f(x)=a\sqrt{x}$ проходит через точку $A(4;6)$. Найдите $f(16)$.
\item Решите уравнение $\sqrt{13-3x}=2$.
\end{enumerate}
\topic{Средний уровень}
\begin{enumerate}[label=\arabic*.,start=5]
\item Упростите при $a>0$: $\displaystyle\left(\frac{a^{3/2}}{a^{-1/2}}\right)^3$.
\item Решите уравнение $\sqrt{4+3x}=x$.
\item Решите уравнение $\sqrt{-12-7x}=-x$.
\end{enumerate}
\topic{Повышенный уровень}
\begin{enumerate}[label=\arabic*.,start=8]
\item Сократите дробь, сохранив ОДЗ: $\displaystyle\frac{a^{1/2}-4a^{1/4}}{a^{1/2}-16}$.
\item Упростите при $a,b>0$: $\displaystyle\left(b^{1\frac15}\cdot\sqrt{a^3}:a^{-\frac3{10}}\right)^{\frac53}$.
\item Решите уравнение $\sqrt{3x^2-4x-7}=x+1$.
\end{enumerate}
\clearpage
\section*{Ответы}
\begin{center}
\renewcommand{\arraystretch}{1.55}
\begin{tabularx}{\linewidth}{@{}lX@{}}
\toprule
№ & Ответ \\
\midrule
1 & $(-\infty;5]$ \\
2 & $(0;6)$ \\
3 & $12$ \\
4 & $3$ \\
5 & $a^6$, при $a>0$ \\
6 & $4$ \\
7 & $-4; -3$ \\
8 & $\dfrac{a^{1/4}}{a^{1/4}+4}$, при $a\ge0$, $a\ne256$ \\
9 & $a^3b^2$, при $a,b>0$ \\
10 & $-1; 4$ \\
\bottomrule
\end{tabularx}
\end{center}
\vspace{8mm}
\begin{center}{\large\color{darkaccent} Итоговая памятка}\end{center}
\checkitem Я различаю ограничения для корня и корня в знаменателе.\par
\checkitem Я правильно отмечаю границы на числовой прямой.\par
\checkitem Я умею восстановить коэффициент по точке графика.\par
\checkitem Я сохраняю минусы и скобки в показателях степеней.\par
\checkitem Я раскладываю дробь на множители и сохраняю ОДЗ.\par
\checkitem Я проверяю знак всей правой части перед возведением в квадрат.\par
\checkitem Я учитываю возможность нулевой правой части.\par
\end{document}
